
\begin{frame}{Prova de Equivalência}
  Lembrando que:

  \begin{minipage}[t]{0.45\textwidth}
    Propriedade Universal do $\Pi$:
    \vspace{0.2cm}
    \[b = \Pi a \iff b = (1 : a \cdot b)\]
  \end{minipage}
  \pause
  \hfill
  \vrule
  \hfill
  \begin{minipage}[t]{0.45\textwidth}
    \begin{calculation}[\calciff]
      I - \Sigma p = \Pi (I - p)
      \step{$a = (I - p)$ e $b = (I - \Sigma p)$}
      I - \Sigma p = I : (I - p) \cdot (I - \Sigma p)
    \end{calculation}
  \end{minipage}
  \vfill
  \pause
  \vfill
  Ou seja, para provar que $I - \Sigma p = \Pi (I - p)$, basta provar que $I - \Sigma p = I : (I - p) \cdot (I - \Sigma p)$.
  \vfill
\end{frame}

\begin{frame}{Prova de Equivalência}
  \begin{calculation}[=]
    I - \Sigma p \pause
    \step{Definição ($I$)} \pause
    (I : I) - \Sigma p \pause
    \step{Definição ($\Sigma$)} \pause
    (I : I) - (0 : p + \Sigma p) \pause
    \step{Definição ($-$)} \pause
    (I - 0) : (I - (p + \Sigma p)) \pause
    \step{Simplificação} \pause
    I : (I - p - \Sigma p)
  \end{calculation}
\end{frame}

\begin{frame}{Prova de Equivalência}
  \begin{calculation}[=]
    I : (I - p - \Sigma p)
    \step{?}
    \vphantom{a}
    \step{?}
    I : (I - p) \cdot (I - \Sigma p)
  \end{calculation}
\end{frame}

\begin{frame}{Prova de Equivalência}
  \begin{calculation}[=]
    I : (I - p - \Sigma p)
    \step{?}
    \vphantom{a}
    \step{Distributividade}
    I : (I - p) \cdot (I - \Sigma p)
  \end{calculation}
\end{frame}

\begin{frame}{Prova de Equivalência}
  \begin{calculation}[=]
    I : (I - p - \Sigma p)
    \step{?}
    I : (I - p - \Sigma p + p \Sigma p)
    \step{Distributividade}
    I : (I - p) \cdot (I - \Sigma p)
  \end{calculation}
\end{frame}

\begin{frame}{Prova de Equivalência}

  \vspace{0.5cm}

  \begin{calculation}[=]
    I : (I - p - \Sigma p)
    \step{$p \Sigma p = 0$}
    I : (I - p - \Sigma p + p \Sigma p)
    \step{Distributividade}
    I : (I - p) \cdot (I - \Sigma p)
  \end{calculation}

  \pause

  \begin{center}
    Se $p \Sigma p = 0$, provamos que $I - \Sigma p = \Pi (I - p)$.
  \end{center}
\end{frame}

\begin{frame}{Prova de Equivalência}
  \begin{itemize}
    \item $p \Sigma p = 0$ diz respeito à corretude do algoritmo (lembrando que $p = qq\t$); \\[0.5cm] \pause
    \item Sabemos provar que $p \Sigma p = 0$ ao mostrar que \[p \Sigma p = \alpha \left( \Sigma p - \left( \Sigma p \right)^2 \right),\] sendo $\alpha$ uma constante, e mostrando que $\Sigma p - \left( \Sigma p \right)^2$ é igual a uma outra stream que é solução de uma recorrência que apresenta a stream nula como única solução (A prova está no relatório).
  \end{itemize}
\end{frame}